\SetRegion{执行摘要}{teal}{compass}
\phantomsection\label{ch1-executive}\pdfbookmark[0]{执行摘要}{ch1-bm-executive}
\kicker{EXECUTIVE SUMMARY / 主要判断}
\PageTitle{参与持续扩张，\protect\\开放能力仍需共同建设}{从生态证据，走向有对象、有条件的行动}
人工智能使开放协作的对象从代码扩展到模型、数据与工具。更多资源能够被取得，并不意味着更多人已经能够用好它们。《2026全球开源发展报告》沿参与、供给、复用、治理与公共价值五段主线，观察开放生态的规模变化、结构差异和持续发展条件。
\par\vspace{1mm}\noindent\includegraphics[width=\linewidth]{figures/new/summary_metrics.pdf}\par\vspace{1mm}
\topic{参与基础与当期活跃增长不同步}
第二章同源记录显示，2022—2026报告年度，开发者总量增至期初的2.32倍，活跃人数增至1.87倍。两者都在扩大，但步调不同。当期活跃不是同一批人的留存率；后续支持需要分别关注进入机会、实际任务与维护接续。\SourceCite{https://github.com/X-lab2017/2026-Global-Open-Source-Development-Report/blob/c51e3441361407b68f10b339d97eff288319b503/preview/content/vmap/source/README.md}{14}
\topic{同一个生态，在不同指标下呈现不同结构}
参与人数、协作贡献和项目影响各自回答不同问题。开发者所在地与项目归属也并不相同。地域规模能够帮助定位资源，贡献关系揭示共同建设，任何单一排名都不足以概括整个生态的发展质量。
\topic{技术供给应落实到分工、接口与维护}
第四章完整全景显示，入口、核心工具与工程支撑承担不同工作；仓库边界与系统功能并不重合。4.6保留PyTorch生态样本全景，不代表正式托管名单。榜单和分类帮助确定调查范围，采用判断仍需具体版本、运行环境、兼容与交付证据。不同技术样本不合成为一个全球总量。
\topic{把资源发布推进到持续使用和独立验证}
第五章将模型供给、声明复用与开放材料分开观察：可下载不等于完整开源，许可开放也不等于材料齐备或运行已经复现。第六章据此提出持续参与、共同维护与结果核查的行动方向；投入应连接具体对象、责任和检验方式，不按热度直接判断公共成效。
\vfill\src{报告主统计期：2025年7月1日至2026年6月30日。第二章自然年快照、第四章异窗样本及第五章期后记录分别保留。各章证据分别解释，不合成总榜或公共成效认证。}
\clearpage
\kicker{EXECUTIVE SUMMARY / 证据与行动}
\PageTitle{四种视角，连接三类共同任务}{从看见生态，到确定谁来做什么}
地图、榜单、技术全景与模型关系分别提供一种证据。只有保留各自的对象和时期，才能把差异转化为可继续调查的问题。下面的行动方向来自跨章分析，仍需在具体场景中检验效果。

\begingroup\sffamily\fontsize{10}{15}\selectfont\renewcommand{\arraystretch}{1.45}
\noindent\begin{tabularx}{\linewidth}{@{}p{28mm}X p{37mm}@{}}\toprule
研究视角 & 本版可以支持的观察 & 下一步要补什么证据\\\midrule
GOSO 生态地图 & 参与基础与当期活跃增长不同步；贡献来源与项目归属分别记录。 & 固定参与批次、持续贡献和退出原因。\\
GOSR 主体排行 & 仓库、项目、机构与模型版本的边界，会改变榜单的含义。 & 实际依赖、版本交付与维护责任。\\
GOSL 技术全景 & 完整清单呈现功能分工；固定样本比较变化，组织案例辨认责任。 & 接口兼容、任务记录与跨版本维护。\\
GOSM 模型生态 & 声明复用、许可条件和材料完整度揭示不同的继续研发条件。 & 具体版本的权利、运行验证与使用结果。\\\bottomrule
\end{tabularx}\endgroup
\src{证据入口：第二章2.1—2.3；第三章3.4—3.6；第四章4.1、4.4、4.6；第五章5.2—5.3。四类证据不共享总体分母，表中不建立因果关系。}

\topic{支持持续参与：让进入机会连接到真实任务}
政府、国际组织、高校与社区共同提供导师、工具和参与条件；按固定批次观察后续贡献与角色成长。年度活跃数用于判断参与规模，不能代替留存或培养效果。对应第六章建议02、07、09。

\topic{维护共同基础：把支持落实到具体版本与责任}
企业、基金会、公共采用方和维护团队核对实际依赖，约定测试、文档、适配与交接工作。基金会支持和项目技术决策分别查证；较高排名本身不能决定资助优先级。对应建议01、04、08。

\topic{检验使用结果：分开记录材料、运行与成效}
发布者说明可获得的材料与权利，实施方记录运行条件和成本，研究与评价人员核验结果及限制。已有技术观察尚缺完整的成本、本地化、维护和结果证据。对应建议03、05、06。

\vfill\src{第六章说明九条建议的对象、交付与约束。附录A说明指标、数据来源与统计口径；开放条件和公共价值分别在相关章节讨论。建议不构成已经验证有效的干预结论。}
\clearpage\SetRegion{执行摘要 · 发展目标}{teal}{compass}
\phantomsection\label{ch1-dimensionoverview}\pdfbookmark[1]{开源发展目标：十个观察维度}{ch1-bm-ten-dimensions}
\kicker{OSDGs / 开源发展目标}
\PageTitle{四个目标，十个观察维度}{让共同方向成为可以持续追问的问题}
\noindent OSDGs提出发展方向，OSDI以十个二级维度组织观察。它们共同关注人才、供给、连接与治理，不以一个总分替代生态的不同侧面。
\begin{center}\includegraphics[width=\linewidth,height=185mm,keepaspectratio]{figures/new/dimension_board.pdf}\end{center}
\src{维度名称与层级沿用第一章。D01—D10为视觉编号；候选指标的定义、数据状态与方法见1.3及附录A。}
\clearpage\SetRegion{执行摘要 · 生态证据}{blue}{globe}
\kicker{EVIDENCE AT A GLANCE / 全球观测}
\PageTitle{三种观察，三种分布}{参与、贡献与项目影响不能互相替代}
\noindent 三张地图采用各自明确的对象与时期；右侧读数帮助准确比较相应前五项。空间形态是阅读入口，不是地区综合实力评分。
\begin{center}\includegraphics[width=\linewidth]{figures/new/summary_atlas.pdf}\end{center}
\src{来源：第二章2.1—2.3的固定归档底表。A为报告年度活跃人数；B为2025年贡献快照；C为2025年前30地域项目影响表。三组分母与覆盖不同，不交叉求和或推算变化。\SourceCite{https://github.com/X-lab2017/2026-Global-Open-Source-Development-Report/blob/c51e3441361407b68f10b339d97eff288319b503/preview/content/vmap/source/README.md}{14}\SourceCite{https://github.com/X-lab2017/2026-Global-Open-Source-Development-Report/blob/c51e3441361407b68f10b339d97eff288319b503/preview/content/cmap/countries.csv}{15}\SourceCite{https://github.com/X-lab2017/2026-Global-Open-Source-Development-Report/blob/c51e3441361407b68f10b339d97eff288319b503/preview/content/imap/influence.json}{16}}



